\documentclass[9pt]{article}

\usepackage[margin=0.2in]{geometry}
\usepackage[T1]{fontenc}
\usepackage[utf8]{inputenc}
\usepackage{lmodern}
\usepackage{listings}
\usepackage{xcolor}
\usepackage{hyperref}

\setlength{\parindent}{0pt}  % <-- removes all paragraph indentation

% Simple JavaScript definition for listings (enough for our examples)
\lstdefinelanguage{JavaScript}{
  keywords={break, case, catch, continue, debugger, default, delete, do, else,
    finally, for, function, if, in, instanceof, new, return, switch, this, throw,
    try, typeof, var, void, while, with, const, let},
  sensitive=true,
  morecomment=[l]{//},
  morecomment=[s]{/*}{*/},
  morestring=[b]",
  morestring=[b]'
}

% Highlight 'fields' and 'allFields' differently inside code
\lstset{
  basicstyle=\ttfamily\small,
  breaklines=true,
  columns=fullflexible,
  emph={fields},                      % group 1
  emphstyle=\color{blue}\bfseries,
  emph={[2]allFields},                % group 2
  emphstyle={[2]\color{red}\bfseries}
}

\begin{document}

\noindent
\textbf{Database Management Systems}%
\hfill%
\textit{Understanding an Aggregation Pipeline for Extracting Field Names in MongoDB}\\
Beata Kubiak, Instructor

\medskip

\section*{Aggregation Pipeline for Listing Field Names}

\noindent
We use the following aggregation pipeline:
\begin{lstlisting}[language=JavaScript]
db.movies.aggregate([
  { $project: { fields: { $objectToArray: "$$ROOT" } } },
  { $unwind: "$fields" },
  { $group: { _id: null, allFields: { $addToSet: "$fields.k" } } }
])
\end{lstlisting}

\subsection*{Step 1: \texttt{\$project} with \texttt{\$objectToArray}}

We use
\begin{lstlisting}[language=JavaScript]
$project: { fields: { $objectToArray: "$$ROOT" } }
\end{lstlisting}

This turns a document like
\begin{lstlisting}[language=JavaScript]
{"_id": ...,
  "title": "...",
  "year": 1999}
\end{lstlisting}

into
\begin{lstlisting}[language=JavaScript]
{"fields": [
    { "k": "_id",   "v": ... },
    { "k": "title", "v": "..." },
    { "k": "year",  "v": 1999 }]}
\end{lstlisting}

So the whole document \texttt{\$\$ROOT} is converted into an array of
key--value pairs under the new field \texttt{\textcolor{blue}{fields}}.

\subsection*{Step 2: \texttt{\$unwind} on \texttt{"\$fields"}}

Next we apply
\begin{lstlisting}[language=JavaScript]
$unwind: "$fields"
\end{lstlisting}

This takes the array in \texttt{\textcolor{blue}{fields}} and turns it into multiple documents:

\begin{lstlisting}[language=JavaScript]
{ fields: { k: "_id",   v: ...   } }
{ fields: { k: "title", v: "..." } }
{ fields: { k: "year",  v: 1999  } }
\end{lstlisting}

Each key--value pair becomes its own document with a single field
\texttt{\textcolor{blue}{fields}} holding an object \texttt{\{ k: ..., v: ... \}}.

\subsection*{Step 3: \texttt{\$group} to collect all field names}

Finally we use
\begin{lstlisting}[language=JavaScript]
$group: { _id: null, allFields: { $addToSet: "$fields.k" } }
\end{lstlisting}

Explanation:
\begin{itemize}
  \item \texttt{\_id: null} means ``put all documents into a single group''. We do not
        want separate groups; we want one combined result.
  \item \texttt{\$addToSet: "\$fields.k"} takes the value of \texttt{fields.k}
        (the field name) from each document and adds it to the set
        \texttt{\textcolor{red}{allFields}}. Using a \emph{set} automatically removes duplicates.
  \item \texttt{\$addToSet} eliminates duplicates, but it only works as part of a \texttt{\$group} stage, which is what gathers values from multiple documents.
\end{itemize}

The result is a single document of the form
\begin{lstlisting}[language=JavaScript]
{
  "_id": null,
  allFields: [
    "_id",
    "title",
    "year",
    // ... all other distinct field names ...
  ]
}
\end{lstlisting}

This gives us the list of every field name that appears in any document
in the collection.

\section*{Which words are ``keywords'' here?}

In the context of this aggregation pipeline:

\begin{itemize}
  \item \textbf{MongoDB aggregation operators (``Mongo keywords''):}
    \texttt{\$project}, \texttt{\$unwind}, \texttt{\$group},
    \texttt{\$objectToArray}, \texttt{\$addToSet}.
    These all start with \texttt{\$} and have special meaning to MongoDB.
  \item \textbf{System variable:} \texttt{\$\$ROOT} is built-in and refers to the current document.
  \item \textbf{User-chosen names:}
    \texttt{\textcolor{blue}{fields}} and
    \texttt{\textcolor{red}{allFields}} are created by us; we could rename them.
\end{itemize}

\end{document}
